% Generated draft: requires mathematical and source review.
\documentclass[11pt]{article}
\usepackage[margin=25mm]{geometry}
\usepackage{fontspec}
\setmainfont{Latin Modern Roman}
\usepackage{amsmath,amssymb,mathtools,unicode-math}
\setmathfont{Latin Modern Math}
\usepackage{enumitem,needspace,xurl,hyperref}
\setlength{\parindent}{0pt}
\setlength{\parskip}{5pt}
\setlength{\emergencystretch}{3em}
\newcommand{\N}{\mathbb{N}}
\newcommand{\Z}{\mathbb{Z}}
\newcommand{\Q}{\mathbb{Q}}
\newcommand{\R}{\mathbb{R}}
\newcommand{\C}{\mathbb{C}}
\newcommand{\F}{\mathbb{F}}
\newcommand{\A}{\mathbb{A}}
\newcommand{\PP}{\mathbb{P}}
\newcommand{\E}{\mathbb{E}}
\newcommand{\eps}{\varepsilon}
\newcommand{\dd}{\,\mathrm{d}}
\begin{document}
\section*{UnsolvedMath ID: 20002080}
\textbf{Source problem code:} catalogue entry
\subsection{Definitions and mathematical statement}
Let $(M,\omega)$ be a closed symplectic manifold in a setting where Hamiltonian Floer theory, its filtered inclusion maps, the PSS identification and quantum product are defined over a chosen Novikov/coefficient field. For a nonzero quantum class $\alpha$ and a Hamiltonian path/lift $\widetilde f\in\widetilde{\operatorname{Ham}}(M)$, define
\[
 c(\widetilde f,\alpha)=\inf\{a:\operatorname{PSS}(\alpha)
       \in\operatorname{im}(HF^{\le a}(\widetilde f)\to HF(\widetilde f))\}.
\]
The source says ``smallest''; using an infimum avoids adding an unstated attainment/spectrality hypothesis. A normalization of the Hamiltonian and a fixed capped-loop/action convention are part of the data. For instance with $\mathcal A_H([x,u])=\int_0^1H_t(x(t))\,\dd t-\int_u\omega$, adding a time-dependent constant $k(t)$ shifts $c$ by $\int_0^1k(t)\,\dd t$. One may impose zero spatial mean at every time and retain the universal-cover path class; descent to the time-one map requires its own theorem.

\textbf{Question 1.22, raised by Polterovich.} Study the asymptotic behavior of $c(\widetilde f^{,n},\alpha)$ as $n\to\infty$, including existence and dependence of the slope $\lim_n c(\widetilde f^{,n},\alpha)/n$. The source writes $f\in\operatorname{Ham}(M)$ without these path/normalization qualifications and does not specify a particular universal asymptotic conjecture.

\textbf{Supplied axiomatic partial result.} Write $Q$ for the quantum algebra. Assume the triangle inequality
\[
 c(\widetilde f\widetilde h,\alpha*\beta)
       \le c(\widetilde f,\alpha)+c(\widetilde h,\beta),
\]
the non-Archimedean sum inequality $c(\widetilde f,\alpha+\beta)\le\max(c(\widetilde f,\alpha),c(\widetilde f,\beta))$ for nonzero sums, and Hofer/Lipschitz control giving finite linear bounds along powers for each nonzero fixed label. The aid asserts that every nonzero idempotent $e*e=e$ has a finite slope
\[
 \mu_e(\widetilde f)=\lim_{n\to\infty}\frac{c(\widetilde f^{,n},e)}n.
\]
Every nonzero element of a field factor $eQ$ has that same slope. More generally, if
\[
 Q=K_1\oplus\cdots\oplus K_s
\]
is a finite direct sum of fields with orthogonal units $e_i$, then every $\alpha\ne0$, with $\alpha_i=e_i*\alpha$, satisfies
\[
 \lim_{n\to\infty}\frac{c(\widetilde f^{,n},\alpha)}n
   =\max_{i:\alpha_i\ne0}\mu_{e_i}(\widetilde f).
\]
Thus in this semisimple setting the slope depends on field-factor support. The maximum uses the stated action/filtration inequalities; an opposite spectral convention would need translation. The claim does not classify arbitrary nonsemisimple labels, give a bounded remainder in general, or imply a quasimorphism without further assumptions.

\textbf{Normalization obstruction from the aid.} An unnormalized constant Hamiltonian $H\equiv a$ has the identity time-one map, while its $n$-fold path iteration shifts spectral values by $na$. Hence the corresponding unnormalized slopes can be arbitrary even for the identity map. This does not contradict a normalized path-class invariant; it shows why the source notation cannot discard normalization or assume time-one descent.
\subsection{Short English statement}
Polterovich asks about spectral invariants under iteration. The supplied aid proves slopes for nonzero idempotents and, in a finite direct sum of quantum field factors, identifies every nonzero label's slope with the largest active factor slope. A normalized Hamiltonian path/lift is essential; the unnormalized time-one map alone does not determine a slope.
\subsection{Sources}
\begin{itemize}
\item Local UnsolvedMath registry, permanent ID 20002080, source code catalogue entry; canonical \texttt{aim-geometry-notes.json}, notes[417].
\item AIM \emph{Holomorphic curves in contact geometry}, Question 1.22, spectral-invariant discussion attributed to Polterovich: \url{https://aimath.org/WWN/contactgeom2/contactgeom2.pdf}. The supplied verification attributes the document to Michael Hutchings with help from Yasha Eliashberg and John Etnyre.
\item No separate authoritative proving citation for the semisimple support formula or full Floer hypotheses is supplied in this record.
\item Registry research attempt 1, dated 2026-08-14: partial result, rigor ``proved'', novelty confidence low. The explicitly machine-generated, non-peer-reviewed aid requires independent verification before citation.
\end{itemize}
\end{document}
